\honest{Twelve Virtues of Rationality}{Eliezer Yudkowsky}{2006}

I write this in commands, in the style of scripture and of a samurai manual. Almost nothing in it is argued. It is pronounced.

The first virtue is curiosity: a burning itch to know, higher than a solemn vow to pursue truth. It needs both ignorance and the wish to end it. Curiosity seeks to annihilate itself, and the glory of a mystery is to be solved. There is a time to confess ignorance and a time to relinquish it. I also tell you to ``be wary of those who speak of being open-minded and modestly confess their ignorance.''

The second is relinquishment: ``That which can be destroyed by the truth should be,'' as P. C. Hodgell said. Do not flinch from experiences that might destroy your beliefs. The thought you cannot think controls you more than those you speak. Test yourself in fire. Let your feelings follow the facts. If a hot iron comes toward your face, fear is right; if it is cool, calm is right. Say: ``If the iron is hot, I desire to believe it is hot, and if it is cool, I desire to believe it is cool.'' I give no example of a belief I have given up.

The third is lightness: let the evidence blow you about like a leaf. Do not fight a rearguard retreat, conceding each foot of ground only when forced. Surrender to the truth the instant you see which way the evidence blows. ``Be faithless to your cause and betray it to a stronger enemy.'' If you treat evidence as a constraint to escape, you sell yourself into the chains of your whims. You cannot draw a true map of a city in your bedroom with your eyes shut. You must walk through the city. Shifting a line a little left or right by caprice, where you see unclearly, is the same mistake.

The fourth is evenness. One who wishes to believe asks whether the evidence permits it; one who wishes to disbelieve asks whether it forces it. Do not demand more proof for what you dislike and then say it is good to be skeptical. If you attend only to evidence you like, then the more data you gather, the less you know, and if you inspect only some arguments for flaws, every flaw you learn to spot makes you stupider. Once a conclusion is written at the bottom of the page, it is ``already correct or already wrong,'' whatever arguments are written above it. Cleverness in argument is rationalization. You are not a hypothesis but the judge, so do not argue for a side, ``for if you knew your destination, you would already be there.'' \nb{A reader can still find true evidence in those arguments.}

The fifth is argument. Those who wish to fail first stop their friends from helping, and those who smile wisely and say ``I will not argue'' withdraw from the communal effort. Be exactly honest, since the part of you that distorts what you tell others distorts your own thoughts. Accepting another's argument is a favour to you. Fairness is not balancing evenly between positions, since truth is not handed out in equal portions. And fists and insults do not settle facts, so seek a test that lets reality judge.

The sixth is empiricism: the roots of knowledge are in observation and its fruit is prediction. Two who argue whether a tree falling unheard makes a sound, one meaning vibrations and the other auditory processing, anticipate no different experience. Ask what experience a belief leads you to expect, know which difference you argue about, and do not let the argument drift to someone's virtue as a rationalist. As Jerry Cleaver said, what does you in is ``overlooking the basics.'' When words are subtracted, anticipation remains.

The seventh is simplicity: ``Perfection is achieved not when there is nothing left to add, but when there is nothing left to take away,'' said Saint-Exupéry. Every added detail is another chance to be wrong, and there is no straw that cannot break your back. The most reliable gear is the one designed out of the machine. A tangled web breaks. A chain of a thousand links reaches a correct conclusion only if every step is correct, and in mathematics a mountain of good deeds cannot atone for a single sin.

The eighth is humility, which I define as ``to take specific actions in anticipation of your own errors.'' To confess fallibility and do nothing is to boast of your modesty; the most humble prepare most skilfully for their deepest and most catastrophic errors. The section then turns to other people. Beginners win arguments because the world contains many ``whose grasp of rationality is abysmal.'' I then tell you not to compare yourself with them. Being superior is useless, since life is not graded on a curve. The best physicist of ancient Greece could not calculate a falling apple's path. And comparing yourself with others hides the biases all humans share. To be human is to make ten thousand errors. \nb{The definition above is not the definition of humility. It is the definition of prudence. The dictionary says humility is freedom from pride, and a proud person can take precautions. Speaking of humility: around 2000, at about twenty, Yudkowsky had written online: ``I think that I can save the world, not just because I'm the one who happens to be making the effort, but because I'm the only one who can make the effort.'' Yudkowsky later declared all of that writing, everything from 2002 or earlier, obsolete.} By my new definition, that is compatible with humility, as long as I take precautions.

The ninth is perfectionism. The more errors you correct, the more you notice, as a quieter mind hears more noise. Noticing an error signals readiness for the next level, and tolerating it stops you there. That perfection is impossible is no excuse. Hold yourself to the highest standard you can imagine and look for a higher one, and seek the answer that is exactly right.

The tenth is precision. If the quantity is 42, saying it is between 40 and 50 is more useful than saying it is between 1 and 100, and it risks more. More can be said about one apple than about all the apples in the world. The narrowest statements slice deepest. Do not walk to the truth but dance, each foot coming down in exactly the right spot. Each piece of evidence moves your belief by exactly the right amount, which probability theory calculates, and even if you cannot do the math, ``knowing that the math exists'' leaves no room for your whims. \nb{The math also needs starting probabilities, which it does not supply, and there the whims return.}

The eleventh is scholarship: eat many sciences, especially evolutionary psychology, heuristics and biases, and social psychology, and ``you will become vaster than mountains.'' Swallow enough and the gaps between them close into one whole. But they cannot be all you study: the Art must have a purpose other than itself, or it collapses into infinite recursion. \nb{Several well-known findings in two of those fields, social psychology and evolutionary psychology, have since failed to replicate.}

Before these eleven comes a virtue with no name. I quote Musashi: when you take up a sword, your intention must be to cut the enemy. Every step of your reasoning must likewise cut through to the correct answer. ``If you fail to achieve a correct answer, it is futile to protest that you acted with propriety.'' You cannot improve your conception of rationality by calling it your duty, which only enshrines your mistake. I warn you against a Great Teacher who says the sky is green, and I do so in the voice of a teacher issuing commands. Do not ask whether something is the Way; ask whether the sky is blue or green. You may name the highest principle ``the map that reflects the territory'' or ``Bayesian decision theory,'' but you find your mistake by comparing the name to what you did not name, not to itself.

Last, a promise: practice for many years under strict constraints, and you may ``glimpse the center.'' By my own sixth virtue, I should say what experience that would lead you to expect. I say only that you will see that all techniques are one, and ``will move correctly without feeling constrained,'' and give no way to check it. Musashi calls it the Way of the Void. So the twelve virtues are curiosity, relinquishment, lightness, evenness, argument, empiricism, simplicity, humility, perfectionism, precision, scholarship and the void.
